\documentclass[12pt,a4paper]{article}
\usepackage[T2A]{fontenc}
\usepackage[utf8]{inputenc}
\usepackage[russian]{babel}
\usepackage{amsmath,amssymb,amsthm}
\usepackage{geometry}
\geometry{margin=2.5cm}
\usepackage{graphicx}
\usepackage{hyperref}
\hypersetup{colorlinks=true,linkcolor=blue,citecolor=blue,urlcolor=blue}
\usepackage{booktabs}
\usepackage{enumitem}

\newtheorem{theorem}{Теорема}
\newtheorem{lemma}{Лемма}
\newtheorem{proposition}{Утверждение}
\newtheorem{remark}{Замечание}
\newtheorem{definition}{Определение}

\title{\textbf{Единая геометродинамическая модель элементарных частиц
как дефектов упругого эфира.\\
Расширенное изложение с полными выкладками}}
\author{А.~Ю.~Дроздов}
\date{24 октября 2026 г.}

\begin{document}
\maketitle

\begin{abstract}
\noindent
Настоящая работа --- расширенная версия препринта о геометродинамической модели
частиц как дефектов упругого эфира. В отличие от обзорного изложения, здесь
\emph{каждое} структурное утверждение сопровождается полным выводом: размерным
анализом, промежуточными формулами и, где уместно, доказательством в форме теоремы.
Цель --- сделать все результаты воспроизводимыми и проверяемыми читателем
<<с карандашом>>. Границы модели (что выводится, что принимается, что открыто)
прослеживаются по ходу каждого доказательства, а не только в итоговом разделе.
\end{abstract}

\tableofcontents
\newpage

\section{Обозначения и соглашения}
Трёхмерный упругий континуум (<<эфир>>) описывается полем смещений $u_i(\mathbf{x},t)$.
Базис размерностей --- $(L,T,M)$. Греческие индексы $a,b=1,2,3$ --- внутренние
(трансляционные) индексы калибровочной группы $T(3)$; латинские $i,j,k$ ---
пространственные. Свёртка по повторяющимся индексам подразумевается.

\begin{definition}[Основные упругие величины]
Модуль сдвига $\mu$ $[ML^{-1}T^{-2}]$, плотность $\rho$ $[ML^{-3}]$, скорость
поперечных волн $c=\sqrt{\mu/\rho}$ $[LT^{-1}]$, вектор Бюргерса $b$ $[L]$.
\end{definition}

\section{Калибровочная структура и словарь Максвелла}

\subsection{Поле дисторсии и тензор дислокаций}
Пластическая несовместность вводится калибровочным полем $\phi^a_i$. Тензор
плотности дислокаций (кручение Картана) определяется как
\begin{equation}
T^a_{ij} = \partial_i \phi^a_j - \partial_j \phi^a_i .
\label{eq:torsion}
\end{equation}
Интеграл по поверхности, натянутой на контур $\partial S$, даёт вектор Бюргерса:
\begin{equation}
b^a = \oint_{\partial S}\phi^a_i\,dx^i = \int_S T^a_{ij}\,dS^{ij}.
\label{eq:burgers}
\end{equation}

\begin{lemma}[Размерность магнитного заряда]
\label{lem:qm}
В упругом базисе магнитный заряд, определённый как поток плотности дислокаций
$q_m=\int_S T^a\,dS$, имеет размерность длины.
\end{lemma}
\begin{proof}
Из \eqref{eq:torsion} следует $[T^a_{ij}]=[\partial\phi]=[\phi]/L$. Поскольку
$\phi^a_i$ безразмерно (градиент смещения, $[\partial u]=1$), имеем
$[T^a]=L^{-1}$. Тогда $[q_m]=[T^a]\cdot[dS]=L^{-1}\cdot L^{2}=L$. $\qed$
\end{proof}

\begin{remark}
Лемма~\ref{lem:qm} снимает кажущееся противоречие с СИ, где магнитный заряд
измеряется в $\text{А}\cdot\text{м}$: в последовательной упругой теории ампер не
является первичной единицей, и размерность заряда диктуется самой теорией.
\end{remark}

\subsection{Квази-максвелловские уравнения}
Вводя ток дислокаций $I^a_i=\partial_t\phi^a_i-\partial_i v^a$ (где $v^a$ ---
скорость), получаем тождества Бьянки, структурно совпадающие с уравнениями
Максвелла \cite{lazar2010}:
\begin{align}
\partial_i {}^{*}T^{a\,i} &= 0, & \text{(нет истоков у }T^a\text{)}\\
\epsilon_{ijk}\partial_j I^a_k + \partial_t T^a_i &= 0. & \text{(аналог закона Фарадея)}
\end{align}
Отсюда словарь: $\mathbf{B}\leftrightarrow T^a$, $\mathbf{E}\leftrightarrow I^a$.

\subsection{Полный вывод квази-максвелловской системы}
Приведём вывод подробно. Лагранжева плотность линейной динамической теории
дислокаций в калибровке $T(3)$ бьётся на упругую и дислокационную части
\cite{lazar2010}:
\begin{equation}
\mathcal{L} = \tfrac12\rho\,v^a v_a - \tfrac12 C_{ijkl}\beta^{i}_{a}\beta^{j a}
 + \tfrac12 D_{ab}\,I^a_i I^{b\,i} - \tfrac14 A_{ab}\,T^a_{ij}T^{b\,ij},
\label{eq:lagr}
\end{equation}
где $\beta$ --- упругая дисторсия, $C_{ijkl}$ --- тензор упругости, а $D_{ab}$,
$A_{ab}$ --- материальные тензоры откликов (аналоги диэлектрической и магнитной
проницаемостей). Полевые напряжённости --- плотность дислокаций $T^a$
\eqref{eq:torsion} и ток $I^a$; отклики определяются как
\begin{equation}
H_a^{\ ij} = -\frac{\partial\mathcal{L}}{\partial T^a_{ij}} = A_{ab}T^{b\,ij},
\qquad
D_a^{\ i} = \frac{\partial\mathcal{L}}{\partial I^a_i} = D_{ab}I^{b\,i}.
\label{eq:constit}
\end{equation}
Варьирование действия $\delta\!\int\mathcal{L}\,d^4x=0$ по калибровочному
потенциалу $\phi^a$ даёт уравнение движения (аналог закона Ампера--Эрстеда):
\begin{equation}
\partial_j H_a^{\ ij} - \partial_t D_a^{\ i} = J_a^{\ i},
\label{eq:ampere}
\end{equation}
где $J_a$ --- ток источника (плотность импульса дислокаций). Два тождества
Бьянки для $T^a=d\phi^a$ (замкнутость формы, $d^2=0$) дают вторую пару:
\begin{equation}
\partial_{[i}T^a_{jk]}=0,\qquad \partial_t T^a_{ij}+\partial_{[i}I^a_{j]}=0.
\label{eq:bianchi}
\end{equation}
Сопоставляя \eqref{eq:constit}--\eqref{eq:bianchi} с уравнениями Максвелла,
получаем словарь таблицы~\ref{tab:corr}: напряжённости $T^a,I^a$ отвечают
$\mathbf B,\mathbf E$; отклики $H_a,D_a$ --- $\mathbf H,\mathbf D$; а сила
Пича--Кёлера на элемент дислокации
\begin{equation}
f_a^{\ i} = (e_a\!\lrcorner\,I^b)\,p_b^{\ i} + (e_a\!\lrcorner\,T^b)\wedge\Sigma_b
\label{eq:pk}
\end{equation}
играет роль силы Лоренца $q\mathbf E + q\mathbf v\times\mathbf B$. Это и есть
формальное основание всей аналогии: электродинамика частицы-дефекта есть
динамика его дислокационного поля.

\section{Частица: прецессирующий дислокационный диполь}

\subsection{Кинематика особой точки}
Кольцо радиуса $a$ лежит в неподвижной плоскости; особая точка (локализованный
заряд $q_m=\Delta b$) бежит по нему: азимут $\varphi(t)=\varphi_0+\omega t$.
Положение, скорость и ускорение:
\begin{align}
\mathbf{r}(t) &= a(\cos\varphi,\ \sin\varphi,\ 0),\\
\mathbf{v}(t) &= a\omega(-\sin\varphi,\ \cos\varphi,\ 0),\quad |\mathbf{v}|=a\omega,\\
\mathbf{W}(t) &= -a\omega^2(\cos\varphi,\ \sin\varphi,\ 0),\quad |\mathbf{W}|=a\omega^2.
\label{eq:accel}
\end{align}

\subsection{Масса как упругое самодействие}
\begin{proposition}[Размерная замкнутость массы]
Энергия самодействия дислокационного кольца $E=\tfrac12\mu b^2 a\,\ln(8a/r_c)$
(Хирт--Лоте \cite{hirth1982}) после отбрасывания логарифма даёт массу
$m=\mu b^2 a/c^2$ с правильной размерностью $[M]$.
\end{proposition}
\begin{proof}
$[\mu b^2 a]=(ML^{-1}T^{-2})(L^2)(L)=ML^{2}T^{-2}=[\text{энергия}]$. Деление на
$c^2$ $[L^2T^{-2}]$ даёт $[M]$. $\qed$
\end{proof}

\subsection{Диполь из пары колец (противофаза)}
Два кольца в $z=\pm d/2$, особые точки в противофазе ($\varphi_{\text{низ}}=
\varphi_{\text{верх}}+\pi$), заряды $\pm q_m$. Мгновенный дипольный момент
$\mathbf{p}=q_m\mathbf{r}_+ - q_m\mathbf{r}_-$. Прямое вычисление:
\begin{equation}
p_z = q_m\cdot d = \text{const},\qquad
|\mathbf{p}_{xy}| = 2 q_m a\ \ (\text{вращается с }\omega).
\end{equation}
Осевая компонента постоянна (статический диполь), поперечная прецессирует.

\subsection{Устойчивость: теорема о центробежном барьере}
\begin{theorem}[Равновесие пары]
\label{th:stab}
Пусть на пару действует притяжение $U_{\text{прит}}=-A/d$ (противоположные
заряды) и сохраняется момент импульса $J=m a^2\omega$. Тогда эффективная энергия
$E(d)=-A/d+J^2/(2md^2)$ имеет единственный минимум при
$d_{\text{eq}}=J^2/(mA)$, и при $\omega\to0$ (то есть $J\to0$) равновесие
исчезает ($d_{\text{eq}}\to0$).
\end{theorem}
\begin{proof}
$E'(d)=A/d^2-J^2/(md^3)=0 \Rightarrow d_{\text{eq}}=J^2/(mA)$.
Вторая производная $E''(d_{\text{eq}})=-2A/d^3+3J^2/(md^4)$; подставляя
$d_{\text{eq}}$, получаем $E''=A^4 m^3/J^6>0$ --- минимум. При $J\to0$
барьерный член обращается в нуль, остаётся только притяжение, и минимум
смещается в $d\to0$ (схлопывание). $\qed$
\end{proof}

\begin{remark}[Аннигиляция]
Схлопывание при потере прецессии есть рекомбинация Френкеля: интерстициальная
петля ($+\Delta b$) и вакансионная ($-\Delta b$) дают суммарный вектор Бюргерса
$\mathbf{b}_++\mathbf{b}_-=0$ --- идеальную решётку. Таким образом прецессия
(спин) --- механизм, предохраняющий частицу-античастицу от самоаннигиляции.
\end{remark}

\subsection{Спин и статистика: период конфигурации}
\begin{proposition}[Период пары и статистика]
Для пары обычных колец конфигурация возвращается при $\varphi\to\varphi+2\pi$
(бозон). Для пары мёбиусовых колец, где знак Бюргерса модулируется как
$\cos(\varphi/2)$, знак восстанавливается лишь при $\varphi\to\varphi+4\pi$
(фермион, спин $1/2$).
\end{proposition}
\begin{proof}
Твист $\cos(\varphi/2)$ антипериодичен по $2\pi$:
$\cos((\varphi+2\pi)/2)=\cos(\varphi/2+\pi)=-\cos(\varphi/2)$, и периодичен по
$4\pi$. Следовательно спинорная (двузначная) структура возникает именно для
мёбиусова звена. $\qed$
\end{proof}

\begin{lemma}[Вектор Франка твиста $\cos(\varphi/2)$]
\label{lem:franck}
Твист вектора Бюргерса $b(\varphi)=\cos(\varphi/2)$ эквивалентен повороту
Франка на фиксированный угол $\Omega=\pi$ за один обход $\varphi:0\to2\pi$,
то есть клиновой \emph{полудисклинации}. Индекс Франка $\nu=\Omega/2\pi=1/2$,
что даёт спин $1/2$.
\end{lemma}
\begin{proof}
Фаза твиста есть $\varphi_{\text{tw}}(\varphi)=\varphi/2$. Полный поворот
направления Бюргерса за обход равен
$\Omega=\oint d\varphi_{\text{tw}}=\int_0^{2\pi}\tfrac12\,d\varphi=\pi$.
За $4\pi$ поворот составляет $2\pi$ (полный возврат), что воспроизводит
спинорную двузначность леммы о периоде. Через спин-связность Осипова
(\texttt{arXiv:1001.4450}) $\Omega_\varphi=i\omega\sigma_3$ с потоком
$\oint\Omega_\varphi\,d\varphi=2\pi\nu$ получаем $\nu=1/2$ и магнитный момент
$\mu=\nu=1/2$ как поток спин-связности (а не ток петли). $\qed$
\end{proof}

\begin{remark}[Мёбиус-диполь и связь с SO(3)-сектором]
Электрон в модели --- это пара мёбиусовых колец с \emph{противоположными}
направлениями вектора Бюргерса: одно кольцо несёт разрыв $+\Delta b$
(полюс~$N$), второе $-\Delta b$ (полюс~$S$); вместе они дают статический
магнитный диполь $q_m d$. Это отличает электрон (мёбиус-пара, период $4\pi$,
фермион) от бозона (пара обычных колец, период $2\pi$). Существенно, что твист
Бюргерса $b(\varphi)=\cos(\varphi/2)$ есть поворот Франка на $\pi$ за обход ---
то есть \emph{дисклинационная} (ротационная, $SO(3)$) степень свободы. Таким
образом мёбиусов сектор модели \emph{уже} лежит в $SO(3)$: магнитный момент
$\mu=\nu=n/2$ выводится как поток спин-связности $\oint\omega_\varphi
\,d\varphi=2\pi\nu$ (Осипов, \texttt{arXiv:1001.4450}), а не как ток петли.
Гидродинамически тот же объект есть вихрь: завихренность $\nabla\times v$
отвечает плотности Франка, циркуляция --- мощности дисклинации. Прецессия
особой точки ($\omega>0$), обязательная для устойчивости пары, и есть этот
вихрь; частица без вращения коллапсирует.
\end{remark}

\section{Магнитный момент: почему диполь, а не ток}

\begin{theorem}[Зануление токового осевого момента]
\label{th:mu0}
Для пары колец в противофазе с противоположными зарядами усреднённый по периоду
осевой магнитный момент <<кольцевого тока>> $\boldsymbol{\mu}_{\text{ток}}=
\tfrac12\oint q\,\mathbf{r}\times\mathbf{v}$ равен нулю.
\end{theorem}
\begin{proof}
Оба кольца бегут в одну сторону, поэтому $\mathbf{r}\times\mathbf{v}$ одинаков
для обоих и направлен вдоль $+z$: $(\mathbf{r}\times\mathbf{v})_z=a^2\omega$.
Но заряды противоположны: $q_{\text{верх}}=+q_m$, $q_{\text{низ}}=-q_m$. Вклады
$\tfrac12 q_m a^2\omega$ и $-\tfrac12 q_m a^2\omega$ сокращаются:
$\mu_{\text{ток},z}=0$. $\qed$
\end{proof}

\begin{remark}
Теорема~\ref{th:mu0} означает, что магнитный момент электрона в модели есть
\emph{статический} дипольный момент $\mu_{\text{маг}}=q_m d$, а не момент
кругового тока. Это согласуется с тем, что <<ток>> по кольцу --- магнитный, и
применение к нему электрической формулы токовой петли некорректно.
\end{remark}

\begin{proposition}[Размерная несостоятельность $g$-фактора по формуле ЭД]
Попытка определить $g$ через $\mu_{\text{маг}}=g(q_m/2m)S$ в упругом базисе даёт
для $g$ размерность $T/L$, а не безразмерную величину; следовательно формула
электродинамики не переносится дословно, и классический $g=2$ в модели колец не
получается (выходит $g=1$).
\end{proposition}
\begin{proof}
$[\mu_{\text{маг}}]=[q_m][d]=L^2$. $[S]=[m a^2\omega]=M L^2 T^{-1}$.
$[(q_m/m)S]=(L/M)(ML^2T^{-1})=L^3T^{-1}$. Тогда
$[g]=[\mu_{\text{маг}}]/[(q_m/m)S]=L^2/(L^3T^{-1})=T/L\neq 1$. $\qed$
\end{proof}

\section{Плёнка Ферми и гравитация: полный вывод}

\subsection{Постановка}
Плёнка Ферми --- трёхмерная граница $\partial\Omega$ четырёхмерного интеграла
действия, натянутая ортогонально мировой трубке. Её вклад в действие:
\begin{equation}
S_{\partial} = \gamma\int_{\partial\Omega} d\Sigma,
\end{equation}
где $\gamma$ --- модуль упругости плёнки $[ML^{-1}T^{-2}]=[\text{Па}]$.

\subsection{Источник деформации}
\begin{proposition}[Положительность источника]
Усреднённый по периоду источник деформации плёнки от прецессирующего заряда
\begin{equation}
\sigma = \mu b^2 \frac{\langle W^2\rangle}{c^2}
= \mu b^2 \frac{\omega^4 a^2}{c^2}
\end{equation}
строго положителен и не зависит от знака заряда.
\end{proposition}
\begin{proof}
Из \eqref{eq:accel} $|\mathbf{W}|=a\omega^2$, поэтому $\langle W^2\rangle=
a^2\omega^4>0$ (усреднение постоянной величины). Знак заряда входит через $b^2$,
то есть выпадает. $\qed$
\end{proof}

\subsection{Закон обратных квадратов}
\begin{theorem}[$1/R^2$ как следствие функции Грина]
\label{th:invsq}
Если деформация плёнки удовлетворяет уравнению равновесия упругой мембраны
\begin{equation}
\gamma\,\nabla^2\Delta(\mathbf{r})=-\sigma\,\delta^3(\mathbf{r}-\mathbf{r}_0),
\label{eq:poisson}
\end{equation}
то потенциал взаимодействия двух источников убывает как $1/R$, а сила --- как
$1/R^2$.
\end{theorem}
\begin{proof}
Уравнение \eqref{eq:poisson} есть уравнение Пуассона. Его фундаментальное решение
(функция Грина лапласиана в $\mathbb{R}^3$) есть
$G(\mathbf{r})=1/(4\pi|\mathbf{r}|)$, откуда
\begin{equation}
\Delta(\mathbf{r})=\frac{\sigma}{4\pi\gamma\,|\mathbf{r}-\mathbf{r}_0|}.
\end{equation}
Энергия второго источника $\sigma_2$ в поле первого:
$U(R)=\sigma_2\Delta(R)=\sigma_1\sigma_2/(4\pi\gamma R)\propto R^{-1}$.
Сила $F=-dU/dR=\sigma_1\sigma_2/(4\pi\gamma R^2)\propto R^{-2}$.
Показатель $-2$ фиксирован размерностью пространства ($n=3$: $G\sim R^{-(n-2)}$),
а не выбором. $\qed$
\end{proof}

\subsection{Два масштаба жёсткости}
\begin{proposition}[Невозможность единой жёсткости]
Модуль, воспроизводящий массу электрона, и планковский модуль гравитации
несовместимы: их отношение $\sim 10^{84}$.
\end{proposition}
\begin{proof}
Из $m_e c^2=\mu_{\text{эфир}}\,r_e^2\lambda_C$ (при $b=r_e$, $a=\lambda_C$):
$\mu_{\text{эфир}}=m_ec^2/(r_e^2\lambda_C)\approx2{,}7\times10^{28}$~Па.
Планковский модуль $\gamma_P=c^7/(8\pi\hbar G^2)\approx1{,}8\times10^{112}$~Па.
Отношение $\gamma_P/\mu_{\text{эфир}}\approx6{,}9\times10^{83}$. $\qed$
\end{proof}

\begin{remark}[Физический смысл модулей]
Величина $\mu_{\text{эфир}}$ --- это \emph{модуль сдвига вакуума}: жёсткость
эфира по отношению к сдвиговой деформации, с размерностью давления
$[\text{Па}]$. Именно через него в базисе $(L,T,M)$ (без си-ампера)
определяются скорость поперечной волны $c=\sqrt{\mu_{\text{эфир}}/\rho}$, масса
частицы $m=\mu_{\text{эфир}}b^2a/c^2$ и кулоновская сила
$F=\mu_{\text{эфир}}b^2/R\cdot(\text{геометрия})$. Физически он отвечает за
сдвиговые деформации (дислокации), из которых построены заряд и масса, и задаёт
физику на масштабе частицы (порядка комптоновской длины $\lambda_C$). Его
значение $2{,}7\times10^{28}$~Па получено \emph{калибровкой} по электрону
(подстановка $m_e$ и $r_e$ в самоэнергию дислокационного кольца), а не выводом
из первооснов.

Величина $\gamma_P=c^7/(8\pi\hbar G^2)$ --- \emph{планковская жёсткость
вакуума}, второй модуль упругости той же размерности $[\text{Па}]$, отвечающий
за гравитационный отклик (деформацию плёнки Ферми) на планковском масштабе
$\ell_P=\sqrt{\hbar G/c^3}$. Гравитационная постоянная $G$ входит в неё как
\emph{входной} параметр. Колоссальное отношение
$\gamma_P/\mu_{\text{эфир}}\sim10^{84}$ есть прямое выражение слабости
гравитации и составляет главную нерешённую проблему модели (происхождение
иерархии двух масштабов).
\end{remark}

\begin{remark}
Таким образом гравитационная постоянная $G$, входящая в $\gamma_P$, есть
\emph{входной} параметр жёсткости эфира. Мы не претендуем вывести её из $(\hbar,c)$;
доказывается лишь \emph{структура} (теорема~\ref{th:invsq}), а не числовое значение.
Это принципиальное отличие от сфальсифицированных версий, где $G$ пряталась в
жёсткость и <<извлекалась>> обратно (тавтология).
\end{remark}

\begin{remark}[Физический смысл $\gamma_P$ и природа иерархии]
Планковский модуль есть планковская \emph{плотность энергии}, делённая на $8\pi$:
\begin{equation}
\gamma_P=\frac{c^7}{8\pi\hbar G^2}=\frac{1}{8\pi}\,\frac{m_P c^2}{\ell_P^3},
\qquad \ell_P=\sqrt{\hbar G/c^3},\quad m_P=\sqrt{\hbar c/G}
\end{equation}
(численное совпадение $<1\%$). Это не произвольная комбинация, а жёсткость
вакуума на планковском масштабе; множитель $8\pi$ тот же, что в уравнениях
Эйнштейна. Аналогично $\mu_{\text{эфир}}=m_e c^2/(r_e^2\lambda_C)$ есть плотность
энергии на масштабе электрона. Отношение обоих модулей
\begin{equation}
\frac{\gamma_P}{\mu_{\text{эфир}}}\sim\frac{m_P}{m_e}\ \text{(с объёмным фактором)}\sim10^{84}
\end{equation}
сводится к отношению масштабов планк/комптон, то есть к малости $m_e/m_P$. Тем
самым иерархия двух жёсткостей \emph{эквивалентна} стандартной \emph{проблеме
иерархии} ($m_e\ll m_P$): модель не решает её, но и не вводит нового произвольного
параметра --- она наследует общеизвестную малость $m_e/m_P$.
\end{remark}

\section{Шесть каналов: единая таблица с выводом законов}
\label{sec:channels}
Каждый закон спадания --- следствие доказанной структуры:
\begin{itemize}[nosep]
\item \textbf{Кулон} $1/R^2$: произведение зарядов-Бюргерсов, прямое кулоновское
ядро (теорема~\ref{th:invsq} для $T^a$-поля).
\item \textbf{Магнитное} $1/R^4$: взаимодействие двух статических диполей
$q_m d$ (производная диполь-дипольного потенциала $\sim1/R^3$).
\item \textbf{Гравитация} $1/R^2$: теорема~\ref{th:invsq} для источника $\sigma>0$.
\item \textbf{Обмен} $1/R^3$: фазовая зависимость $\cos2\Delta\varphi$,
короткодействие.
\item \textbf{Слабое} $e^{-R/R_W}$: массивный переносчик, киральность твиста
(нарушение чётности).
\item \textbf{Сильное} $\propto R$: частичный вектор Бюргерса $b/3$ и лента
дефекта упаковки (линейный потенциал, конфайнмент).
\end{itemize}
Два безразмерных числа замыкают систему:
\begin{equation}
\alpha=\frac{ke^2}{\hbar c}=\frac{1}{137{,}04},\qquad
\alpha_G=\Big(\frac{m_e}{m_{\text{Planck}}}\Big)^2=1{,}75\times10^{-45},\qquad
\frac{\alpha_G}{\alpha}=2{,}4\times10^{-43}.
\end{equation}

\section{Фальсифицируемость}
Модель даёт проверяемые предсказания. Для магнитно-гравитационных устройств
(тип Рощина--Година) относительная поправка к весу от прецессии оценивается как
\begin{equation}
\frac{\Delta G}{G}\sim\frac{W R}{c^2}\sim 10^{-13},
\end{equation}
что на $\sim\!12$ порядков меньше заявленных эффектов ($\sim\!35\%$); последние,
если реальны, \emph{опровергали} бы модель, а не подтверждали. Это --- признак
научной, а не подгоночной теории.

\section{Границы: доказано / принято / открыто}
\paragraph{Доказано (с выводом выше).} Размерность заряда (лемма~\ref{lem:qm});
замкнутость массы; устойчивость и аннигиляция (теорема~\ref{th:stab});
статистика из периода; зануление токового момента (теорема~\ref{th:mu0});
положительность источника гравитации; закон $1/R^2$ (теорема~\ref{th:invsq});
необходимость двух масштабов.

\paragraph{Принято.} Значения $\alpha$, $\alpha_G$; $G$ в составе $\gamma_P$;
два масштаба жёсткости.

\paragraph{Открыто.} Численные значения $\alpha,\alpha_G$; строгий $g=2$
(в классике $g=1$; $g=2$ --- только на спинорном уровне \cite{osipov2010});
динамическая жёсткость плёнки (размерный множитель в источнике); константы
слабого и сильного секторов.

\paragraph{Тензорный ($J{=}2$) канал гравитации --- диагноз.} Использованный
источник $\sigma\propto\langle W^2\rangle$ есть \emph{след} (скаляр, $J=0$)
квадрупольного тензора $\langle W_iW_j\rangle$. Настоящий тензорный гравитон
(спин-2 по Хелю \cite{hehl1995}) живёт в бесследовой части
$Q_{ij}=\langle W_iW_j\rangle-\tfrac13\delta_{ij}\langle W^2\rangle$. Для
одиночного прецессирующего кольца с осью $\mathbf n$ прямой расчёт даёт
$\langle W_iW_j\rangle=\tfrac12\omega^4a^2(\delta_{ij}-n_in_j)$, то есть
$Q_{ij}\neq0$ (сплюснутый квадруполь в плоскости прецессии, собственные
значения $\{-\tfrac13,+\tfrac16,+\tfrac16\}\omega^4a^2$). Однако усреднение по
изотропно ориентированным осям обращает $\langle Q_{ij}\rangle\to0$: для
неполяризованного ансамбля выживает только скаляр $\langle W^2\rangle$. Таким
образом «скалярность» гравитации модели --- следствие ориентационного
усреднения, а не отсутствия $J=2$. Тензорный канал проявлялся бы лишь для
\emph{поляризованной} (спин-выровненной) материи. Это диагноз и проверяемое
предсказание, а не вывод тензорного гравитона: явное уравнение поля для
$Q_{ij}$ остаётся открытым.

\section{Заключение}
Приведён полный вывод всех структурных результатов модели. Доказательства
элементарны (размерный анализ, функция Грина, теорема вириала), но именно их
отсутствие в прежних версиях приводило к подгонкам. Честно разделены доказанное,
принятое и открытое. Следующие шаги --- спинорный вывод $g=2$ и англоязычная
версия для архива.

\appendix
\section{Численное приложение: графики}
Расчёты выполнены в модельных упругих единицах ($a=\omega=b=c=1$). Абсолютные
значения условны; физический смысл несут \emph{законы спадания} и
\emph{соотношения}, доказанные в основном тексте.

\begin{figure}[h]
\centering
\includegraphics[width=0.75\linewidth]{figures_unified/channel_scaling.png}
\caption{Законы спадания каналов (лог-лог): кулон и гравитация $\sim1/R^2$,
магнитное $\sim1/R^4$, обмен $\sim1/R^3$, слабое --- экспоненциальное
короткодействие. Наклоны отвечают теореме~\ref{th:invsq} и диполь-дипольной
структуре.}
\end{figure}

\begin{figure}[h]
\centering
\includegraphics[width=0.75\linewidth]{figures_unified/stability.png}
\caption{Устойчивость пары (теорема~\ref{th:stab}): $E(d)=-A/d+J^2/(2md^2)$.
Минимум при $d_{\text{eq}}=J^2/(mA)$ растёт с моментом $J$; при $J\to0$
барьер исчезает (аннигиляция).}
\end{figure}

\begin{figure}[h]
\centering
\includegraphics[width=0.75\linewidth]{figures_unified/twist_period.png}
\caption{Твист вектора Бюргерса: $\cos(\varphi/2)$ антипериодичен по $2\pi$ и
возвращается лишь за $4\pi$ (спинор, фермион), в отличие от обычного
$\cos\varphi$ с периодом $2\pi$ (бозон).}
\end{figure}

\begin{thebibliography}{99}
\bibitem{fermi1923} E.~Fermi, Nuovo Cimento \textbf{25}, 159 (1923).
\bibitem{lazar2010} M.~Lazar, Math.\ Mech.\ Solids \textbf{16}(3), 253 (2010).
\bibitem{kadic1983} A.~Kadi\'c, D.~G.~B.~Edelen, Lecture Notes in Physics \textbf{174}, Springer (1983).
\bibitem{osipov2010} E.~A.~Kochetov, V.~A.~Osipov, arXiv:1001.4450 (2010).
\bibitem{hehl1995} F.~W.~Hehl et al., Phys.\ Rep.\ \textbf{258}, 1 (1995).
\bibitem{hirth1982} J.~P.~Hirth, J.~Lothe, \emph{Theory of Dislocations}, 2nd ed., Wiley (1982).
\bibitem{mistura1995} L.~Mistura, Int.\ J.\ Eng.\ Sci.\ \textbf{33}, 2149 (1995).
\bibitem{rohrlich1960} F.~Rohrlich, Am.\ J.\ Phys.\ \textbf{28}, 639 (1960).
\end{thebibliography}

\end{document}
